\documentclass[10pt,a4paper]{amsart}
\usepackage{amsmath,amssymb,amsthm,graphicx}
\numberwithin{equation}{section}

\newtheorem{Theorem}{Theorem}[section]
\newtheorem*{Theorem*}{Theorem}
\newtheorem{Corollary}[Theorem]{Corollary}
\newtheorem{Lemma}[Theorem]{Lemma}
\newtheorem{Proposition}[Theorem]{Proposition}
 { \theoremstyle{definition}
\newtheorem{Definition}[Theorem]{Definition}
\newtheorem{Note}[Theorem]{Note}
\newtheorem{Example}[Theorem]{Example}
\newtheorem{Remark}[Theorem]{Remark} }

\def\Var{\mathrm{Var}}
\def\Nov{\mathrm{Nov}}
\def\Com{\mathrm{Com}}
\def\As{\mathrm{As}}
\def\Lie{\mathrm{Lie}}
\def\Perm{\mathrm{Perm}}
\def\Zinb{\mathrm{Zinb}}
\def\Pois{\mathrm{Pois}}
\def\Jord{\mathrm{Jord}}
\def\SJord{\mathrm{SJord}}
\def\pre{\mathrm{pre}}
\def\RB{\mathrm{RB}}
\def\wt {\mathop {\fam 0 wt}\nolimits}
\def\id {\mathop {\fam 0 id}\nolimits}
\def\Der {\mathop {\fam 0 Der}\nolimits}
\def\End {\mathop {\fam 0 End}\nolimits}
\def\oo#1{\mathbin{{}_{(#1)}}}
\def\oon#1{\mathbin{{\circ}_{(#1)}}}

\def\Red#1{\textcolor{red}{#1}}
\def\Blue#1{\textcolor{blue}{#1}}
\usepackage{cleveref}
\crefname{originallemma}{original lemma}{original lemmas}
\Crefname{originallemma}{Original lemma}{Original lemmas}
\makeatletter
\newlabel{lem:Weight-criterion}{{internal}{source}{}{sourcepointer-1}{}}
\newlabel{lem:Weight-criterion@cref}{{[originallemma][0][]internal}{[1][0][]source}{}{}{}}
\makeatother
\title{Algebras over the derived operad are special}
\author{Kolesnikov, Pavel, Mashurov, Farukh, Sartayev, Bauyrzhan}
\date{doi:10.3842/SIGMA.2024.017; CC BY 4.0}
\begin{document}
\maketitle
\setcounter{section}{3}
\setcounter{Theorem}{4}
\noindent\textit{Editorial note.} All delivered mathematical text is the author's own, copied from the pinned SIGMA article source; no statement or proof step was rewritten, reordered or removed. The counted statement is the labeled theorem environment at source lines 733-738 and the counted proof is the proof environment at source lines 740-759.\par\medskip
\noindent\textit{Editorial note.} The retained preamble of this unit was generated mechanically for this delivery: it copies the author's own package, numbering, macro and theorem declaration block from the pinned source (source lines 9-43), with the publisher class and page setup lines removed.\par\medskip
\noindent\textit{Allowed standard prerequisites (editorial).}\begin{itemize}
\item The embedding of the derived operad into the free operad on the doubled variables. same article, the inclusion used at proof lines 744-746. Allows a presentation of an algebra over the derived operad to be transported into the larger free operad.
\item The weight criterion for membership in the derived operad. same article, the lemma at source lines 562-568. Identifies the elements of the derived operad inside the free operad as the weight homogeneous ones, which is what makes the differential ideal generated by the relations coincide with the relations themselves.
\end{itemize}\medskip
\ifdefined\hypertarget\hypertarget{sourcepointer-1}{}\fi
\noindent\textit{Source reference.} Original lemma internal: sources/sigma-2024-017-on-pre-novikov-algebras-and-derived-zinbiel/pinned-original-source.tex, source lines 562--568. The counted statement or proof refers to that the author's preceding lemma giving the weight criterion for the derived operad; it is the author's own earlier display, definition or statement in the same article and the whole argument of the counted span is retained verbatim.\par\medskip
\begin{Theorem}\label{thm:Weight-Embedding}
If
$\Var $ is a binary operad such that
$\Var\circ \Nov = \Var\otimes \Nov $
then every $D\Var $-algebra is special.
\end{Theorem}

\begin{proof}
Suppose $V$ is a $D\Var $-algebra.
Then $V$ may be presented as a quotient
of a free algebra $D\Var\langle X\rangle $ modulo an ideal $I$.
Consider the embedding $D\Var\langle X\rangle \subset \Var\bigl\langle X^{(\omega )}\bigr\rangle $
and denote by~$J$ the differential ideal of $\Var\bigl\langle X^{(\omega )}\bigr\rangle $
generated by $I$.
Then $U = \Var \bigl\langle X^{(\omega )}\bigr\rangle /J$ is the universal enveloping
differential $\Var $-algebra of $V$. It remains to prove that $J\cap D\Var\langle X\rangle = I$,
namely, the ``$\subseteq $'' part is not a trivial one.

Assume $f\in J$. Then there exists a family of (differential) polynomials
$\Phi_i \in \Var\bigl\langle (X\cup\{t\})^{(\omega )} \bigr\rangle $
such that $f = \sum_i \Phi_i|_{t=g_i}$, for some $g_i\in I$.
If, in addition, $f\in D\Var\langle X\rangle $ then
$\wt (f)=-1$. Since $\wt g_i=-1$, we should have $\wt \Phi_i=-1$ for all $i$.
By Lemma~\ref{lem:Weight-criterion}, every polynomial $\Phi_i$ may be
represented as an element of $D\Var\langle X\cup\{t\}\rangle $,
so $\Phi_i|_{t=g_i} \in I$ for all $i$, and thus $f\in I$.
\end{proof}

\end{document}
